%! AP Statistics — Unit 1, Lesson 1.8 — Guided Notes
\documentclass[10pt]{article}
\usepackage{apstats-article}
\usepackage{apstats-boxes}

%=============================================================================
\begin{document}

\pageheader{Unit 1, Lesson 1.8}{Guided Notes}
\namedateperiod

\vspace{0.10in}

% ── OBJECTIVE ────────────────────────────────────────────────────────────────
\begin{objectivebox}
\small By the end of this lesson, I will be able to\dots\\[4pt]
\writeline\\[1pt]
\writeline\\[1pt]
\writeline
\end{objectivebox}

\vspace{0.10in}

% ── VOCABULARY ───────────────────────────────────────────────────────────────
\begin{vocabbox}
\begin{multicols}{2}
\termblanklong{Boxplot (box-and-whisker plot)}
\termblanklong{Five-number summary}
\termblanklong{Box (IQR box)}
\columnbreak
\termblanklong{Whisker}
\termblanklong{Modified boxplot}
\termblanklong{Fence (lower / upper)}
\end{multicols}
\end{vocabbox}

\vspace{0.10in}

% ── HOOK ─────────────────────────────────────────────────────────────────────
\begin{hookbox}
\small\textbf{Five numbers $\to$ one picture}

\vspace{5pt}
A restaurant chain recorded the delivery times (minutes) for 15 orders one evening.
The five-number summary is: $\min = 12$,\ $Q_1 = 20$,\ $M = 27$,\ $Q_3 = 35$,\ $\max = 48$.

\vspace{4pt}
\begin{center}
\begin{tikzpicture}[scale=1.0]
  \draw[thick] (0.6,0) -- (6.6,0);
  \foreach \v/\lbl in {0.60/10, 1.20/15, 1.80/20, 2.40/25, 3.00/30,
                        3.60/35, 4.20/40, 4.80/45, 5.40/50} {
    \draw (\v,-0.12) -- (\v,0.12);
    \node[below, font=\small] at (\v,-0.12) {\lbl};
  }
  \node[below, font=\small\itshape] at (3.0,-0.56) {Delivery time (minutes)};
  % Five-number: min=12 (x=0.84), Q1=20 (x=1.80), M=27 (x=2.64), Q3=35 (x=3.60), max=48 (x=5.16)
  % scale: (value-10)*0.12 + 0.60
  \draw[thick, navy] (0.84, 0.38) -- (1.80, 0.38);  % left whisker
  \draw[thick, navy] (1.80, 0.10) rectangle (3.60, 0.66);  % box
  \draw[thick, navy] (2.64, 0.10) -- (2.64, 0.66);  % median
  \draw[thick, navy] (3.60, 0.38) -- (5.16, 0.38);  % right whisker
  \draw[thick, navy] (0.84, 0.22) -- (0.84, 0.54);  % left cap
  \draw[thick, navy] (5.16, 0.22) -- (5.16, 0.54);  % right cap
\end{tikzpicture}
\end{center}

\vspace{4pt}
What does each part of this picture represent? We will build the vocabulary to
describe every element.
\end{hookbox}

\vspace{0.08in}

% ── SECTION 1: ANATOMY OF A BOXPLOT ──────────────────────────────────────────
\begin{notesbox}{Anatomy of a Boxplot}
\small
\begin{multicols}{2}

\textbf{\textcolor{navy}{The Five-Number Summary:}}
\begin{itemize}[itemsep=3pt]
  \item $\min = $ \blank{2cm} (left whisker end)
  \item $Q_1 = $ \blank{2cm} (left edge of box)
  \item $M = $ \blank{2cm} (line inside box)
  \item $Q_3 = $ \blank{2cm} (right edge of box)
  \item $\max = $ \blank{2cm} (right whisker end)
\end{itemize}

\vspace{6pt}
\textbf{\textcolor{navy}{What each section tells you:}}
\begin{itemize}[itemsep=3pt]
  \item Each of the 4 sections holds $\approx$ \blank{1.2cm}\% of data.
  \item The box spans the middle \blank{1.2cm}\% of data.
  \item $\text{IQR} = Q_3 - Q_1 = $ \blank{2cm}
\end{itemize}

\columnbreak

\textbf{\textcolor{navy}{Reading values from the hook boxplot:}}

\vspace{4pt}
\renewcommand{\arraystretch}{1.5}
\begin{tabularx}{\columnwidth}{@{} l X @{}}
Range: & $\max - \min = \blank{2cm}$ \\
IQR:   & $Q_3 - Q_1 = \blank{2cm}$ \\
Spread of lower 50\%: & $M - \min = \blank{2cm}$ \\
Spread of upper 50\%: & $\max - M = \blank{2cm}$ \\
\end{tabularx}

\vspace{6pt}
\textit{Skewness hint:} If the right whisker is \blank{2.5cm} than the left,
the distribution is skewed \blank{2cm}.

\end{multicols}
\end{notesbox}

\vspace{0.08in}

% ── SECTION 2: MODIFIED BOXPLOT ───────────────────────────────────────────────
\begin{notesbox}{Modified Boxplot \& the Outlier Rule}
\small
\begin{multicols}{2}

\textbf{\textcolor{navy}{When do we use a modified boxplot?}}\\
Whenever there is at least one \blank{3cm}.

\vspace{6pt}
\textbf{\textcolor{navy}{The $1.5 \times \text{IQR}$ Fence Rule:}}
\begin{itemize}[itemsep=4pt]
  \item Lower fence $= Q_1 - 1.5(\text{IQR})$
  \item Upper fence $= Q_3 + 1.5(\text{IQR})$
  \item Outlier: any value \blank{3cm} a fence
\end{itemize}

\vspace{6pt}
\textbf{\textcolor{navy}{In a modified boxplot:}}
\begin{itemize}[itemsep=3pt]
  \item Whisker extends to the last \blank{2.5cm} value.
  \item The fence is \blank{2.5cm} (not drawn).
  \item Outliers are plotted as \blank{2cm} ($\bullet$).
\end{itemize}

\columnbreak

\textbf{\textcolor{navy}{Example --- Modified Delivery Times:}}\\[4pt]
Five-number summary: $\min=12$, $Q_1=20$, $M=27$, $Q_3=35$, and a
new maximum of $\mathbf{62}$ (the rest of the data is unchanged).
IQR $= 15$.

\vspace{4pt}
Upper fence $= 35 + 1.5(15) = \blank{2.5cm}$\\[3pt]
Is 62 an outlier? \blank{2cm}\\[3pt]
Last non-outlier $= \blank{2cm}$ (whisker ends here)

\vspace{6pt}
\begin{tikzpicture}[scale=0.92]
  % Axis: 10–70, scale (v-10)*0.096 + 0.20; label every 10
  \draw[thick] (0.15,0) -- (6.20,0);
  \foreach \v/\lbl in {0.20/10, 1.16/20, 2.12/30, 3.08/40, 4.04/50, 5.00/60} {
    \draw (\v,-0.10) -- (\v,0.10);
    \node[below, font=\scriptsize] at (\v,-0.10) {\lbl};
  }
  \node[below, font=\scriptsize\itshape] at (3.0,-0.44) {Delivery time (min)};
  % min=12: (12-10)*0.096+0.20=0.392; Q1=20: 1.16; M=27: 1.832; Q3=35: 2.76
  % Last non-outlier=48: (48-10)*0.096+0.20=3.848; outlier 62: 5.552
  \draw[thick, navy] (0.392, 0.32) -- (1.16, 0.32);
  \draw[thick, navy] (1.16, 0.08) rectangle (2.76, 0.56);
  \draw[thick, navy] (1.832, 0.08) -- (1.832, 0.56);
  \draw[thick, navy] (2.76, 0.32) -- (3.848, 0.32);  % whisker to last non-outlier
  \draw[thick, navy] (0.392, 0.18) -- (0.392, 0.46);
  \draw[thick, navy] (3.848, 0.18) -- (3.848, 0.46);
  \filldraw[navy] (5.552, 0.32) circle (0.09);  % outlier dot
\end{tikzpicture}

\end{multicols}
\end{notesbox}

\vspace{0.08in}

% ── SECTION 3: BOXPLOT vs HISTOGRAM ──────────────────────────────────────────
\begin{notesbox}{Boxplot vs.\ Histogram}
\small
\begin{multicols}{2}

\textbf{\textcolor{navy}{Boxplot CAN show:}}
\begin{itemize}[itemsep=3pt]
  \item \blank{3cm} (median)
  \item \blank{3cm} (IQR, range)
  \item Presence of \blank{2.5cm}
  \item Approximate \blank{2cm}
\end{itemize}

\columnbreak

\textbf{\textcolor{navy}{Boxplot CANNOT show:}}
\begin{itemize}[itemsep=3pt]
  \item Exact \blank{3cm} (unimodal? bimodal?)
  \item \blank{3cm} or clusters
  \item Exact \blank{3cm} of individual values
  \item Whether data is \blank{3cm} or discrete
\end{itemize}

\end{multicols}

\vspace{4pt}
\textit{Rule of thumb:} Use a boxplot for \blank{3cm}; use a histogram
for \blank{3.5cm}.
\end{notesbox}

\vspace{0.08in}

% ── REFLECTION ───────────────────────────────────────────────────────────────
\begin{tcolorbox}[colback=goldbg, colframe=goldacc, boxrule=0.8pt, arc=2mm,
    left=4mm, right=4mm, top=2mm, bottom=2mm]
\small\textbf{Quick Check:} A modified boxplot has $Q_1 = 40$, $M = 52$,
$Q_3 = 60$, a left whisker to 30, a right whisker to 78, and an outlier dot at 95.\\[4pt]
IQR $= \blank{2cm}$. \quad
Upper fence $= \blank{3cm}$. \quad
Is 95 above the fence? \blank{2cm}. \quad
Range $= \blank{2cm}$.
\end{tcolorbox}

\end{document}
